\uuid{mVsg}
\chapitre{Statistique}
\niveau{L2}
\module{Analyse de données}
\sousChapitre{Tests d'hypothèses, intervalle de confiance}
\titre{Consommation d'oxygène}
\theme{tests d'hypothèses}
\auteur{}
\datecreate{2022-12-14}
\organisation{AMSCC}
\difficulte{2}
\contenu{
\texte{On s'intéresse à la consommation maximale en oxygène d'un individu dans une population âgée, en ml/kg/min. Pour un groupe de contrôle, il a été montré que les mesures suivent une loi normale d'espérance $25{,}5$ et d'écart-type $\sigma=6$. On pense qu'une population de malades (Parkinson) a des capacités cardio-respiratoires plus limitées. On souhaite tester si, dans ce groupe, l'espérance $\mu$ est plus faible.
On modélise les mesures de $15$ malades par des variables $X_1,\ldots,X_{15}$ indépendantes de loi $\mathcal N(\mu,\sigma)$ ; on admet que l'écart-type vaut aussi $\sigma=6$ chez les malades et qu'il est connu. On prend un risque de première espèce $\alpha=5\%$.}
\begin{enumerate}
\item \texte{L'objectif est de décider entre $H_0:\mu=25{,}5$ (absence d'effet de la maladie) et $H_1:\mu<25{,}5$ (existence de l'effet). On note $\overline X=\frac1{15}\sum_{i=1}^{15}X_i$ la moyenne empirique des mesures.}
\begin{enumerate}
\item \question{La moyenne observée dans l'échantillon est $24{,}0$. Quelle est la décision à prendre à la suite de ce test ?}
\indication{L'écart-type étant connu, la variable de décision suit une loi normale sous $H_0$.}
\reponse{Sous $H_0$, $\overline X\sim\mathcal N\left(25{,}5\,;\frac6{\sqrt{15}}\right)$, avec $\frac6{\sqrt{15}}\simeq1{,}5492$. La variable de décision
\[
Z=\frac{\overline X-25{,}5}{\frac{6}{\sqrt{15}}}
\]
suit donc la loi $\mathcal N(0,1)$. Le test est unilatéral à gauche : on rejette $H_0$ si $Z\leq-u$, avec $u=\Phi^{-1}(0{,}95)\simeq1{,}6449$.
La valeur observée est
\[
z_{\mathrm{obs}}=\frac{24-25{,}5}{1{,}5492}\simeq-0{,}968>-1{,}6449.
\]
On ne rejette pas $H_0$ au risque de $5\%$ ; la $p$-valeur vaut $\Phi(-0{,}968)\simeq0{,}166$. Ce non-rejet ne démontre pas l'absence d'effet : il signifie que l'écart observé reste compatible avec $H_0$ pour cet effectif.}
\item \question{Déterminer la moyenne critique $\mu_c$, c'est-à-dire la valeur observée en deçà de laquelle le test conduit à un rejet de $H_0$.}
\reponse{On rejette $H_0$ lorsque $\overline x\leq\mu_c$, avec
\[
\mu_c=25{,}5-\Phi^{-1}(0{,}95)\frac6{\sqrt{15}}\simeq25{,}5-1{,}6449\times1{,}5492\simeq22{,}95.
\]
La moyenne observée $24{,}0$ est supérieure à $\mu_c$, ce qui confirme le non-rejet.}
\end{enumerate}
\item \texte{La professeure responsable du service pense qu'à partir de la valeur $\mu=23{,}5$, la différence est cliniquement importante. Elle souhaite savoir quel risque elle prend en ne rejetant pas $H_0$ lorsque $\mu=23{,}5$.}
\begin{enumerate}
\item \question{Quelle est la loi de $\overline X$ lorsque $\mu=23{,}5$ ?}
\reponse{Les $X_i$ suivent alors la loi $\mathcal N(23{,}5\,;6)$ et sont indépendantes ; leur moyenne suit la loi
\[
\mathcal N\left(23{,}5\,;\frac6{\sqrt{15}}\right),\qquad\frac6{\sqrt{15}}\simeq1{,}5492.
\]
L'écart-type de la moyenne est divisé par $\sqrt{15}$ par rapport à celui d'une mesure.}
\item \question{En déduire le risque de deuxième espèce et la puissance de ce test pour $\mu=23{,}5$. Commenter le résultat.}
\reponse{On ne rejette pas $H_0$ lorsque $\overline X>\mu_c$. Pour $\mu=23{,}5$,
\[
\beta=\mathbb P(\overline X>\mu_c)
=1-\Phi\left(\frac{22{,}95-23{,}5}{1{,}5492}\right)
=1-\Phi(-0{,}354)\simeq0{,}638.
\]
La puissance vaut $1-\beta\simeq0{,}362$. Avec $15$ malades, le test ne détecte un écart de $2$ ml/kg/min qu'environ une fois sur trois : un non-rejet est donc peu informatif. Pour mieux détecter cet écart, il faudrait augmenter la taille de l'échantillon ou accepter un risque $\alpha$ plus élevé.
\par Au tableur, la moyenne critique s'obtient en B1 par
\par\texttt{=25,5-LOI.NORMALE.STANDARD.INVERSE.N(0,95)*6/RACINE(15)}
\par puis la puissance par \texttt{=LOI.NORMALE.N(B1;23,5;6/RACINE(15);VRAI)}.}
\end{enumerate}
\end{enumerate}
}
